\scanpage{309}
\mspage{(8)}
ou au choix (moyennant \deleted{\unsure{une}} l'équivalence de catégories canonique
\[
 \sigma:(A_*)^{.}\underset{=}{\approx}(A^{.})_*\qquad)
\]
un objet $1/2$ simplicial dans la catégorie $A^{.}$ des complexes de cochaines, \unsure{noté|soit}
\[
 n\longmapsto K_{[n]}{}^{.}.
\]
On interprète de $=$ les objets des autres catégories mixtes \deleted{$\unsure{1/2\text{ simpliciales}}$}.
\[
 \left\{\begin{aligned}
 &A_{*.},\, A_{.*},\, A^*{}_{\,.}{}^{\deleted{\unsure{*}}},\,\deleted{\unsure{A^{.}{}_*}}\\
 &A^{*.}_{\deleted{\unsure{*}}},\, A^{.*}_{\deleted{\unsure{*}}},\, A^{.}{}_*\,\deleted{\unsure{A_*{}^{.}}}.
 \end{aligned}\right.\tag*{(*)}
\]
De façon analogue, les objets $K^{..}$ de $(A^{.})^{.}$ sont les bicomplexes de cochaines, les objets $K_{..}$ de $(A_.)_.=A_{..}$ sont les bicomplexes de chaines, enfin les objets $K^{.}{}_.$ de $(A_.)^{.}$ ou $K_.{}^{.}$ de $(A^{.})_.$ sont les bicomplexes mixtes \inserted{chaines-cochaines} (de chaines pour un indice, de cochaines pour un autre). Dans le cas de $K^{..}$ ou $K_{..}$ les complexes simples \unsure{associé|associés} est un complexe de cochaines resp de chaines, tandis que dans le cas mixte $K^{.}{}_.$ ou $K_.{}^{.}$, on trouve des complexes simples avec des objets en degrés positifs ou négatifs à priori (Il y aurait lieu de fixer une convention de degré, si pour le degré total l'op \unsure{diff|diff.} est de degré $+1$ ou $-1$, suivant
% archive sheet number: 310
